% !TEX root = ../main.tex
% !TEX program = lualatex
\documentclass[../main.tex]{subfiles}
\IfSubfilesClassLoaded{\externaldocument{\subfix{../build/main}}}{}
% =============================================================================
% CHAPTER 27: CIVPERS (CIVILIAN PERSONNEL)
% DESCRIPTION: CIVPERS overview, workforce shaping, \ac{eeo}, employee and labor relations.
% =============================================================================
\AtBeginDocument{\chapter{EDO Study Guide}}{}
\begin{document}
\IfSubfilesClassLoaded{\chapter{EDO Study Guide}}{}
%====================
% Contents here
%====================

\section{Civilian Personnel Overview}

\subsection{Summary}
\begin{description}
  \item[\textbf{Why it matters.}] \ac{civpers} management is a major \ac{edo} responsibility and directly
  impacts readiness across \ac{navsea} and warfare center commands~\autocite{edo-6-1-1-civpers-overview-2025}.
  \item[\textbf{Myths vs.\ reality.}] Common myths about civil service (e.g., poor performance or
  inflexibility) are often inaccurate; effective leaders treat civilians as equal teammates while
  following the governing rules~\autocite{edo-6-1-1-civpers-overview-2025}.
  \item[\textbf{Guidance sources.}] Civilian personnel policy flows from Congress, executive orders,
  \ac{opm}, the \ac{dod}, \ac{secnav}, \ac{navsea} guidance, and negotiated labor agreements~\autocite{edo-6-1-1-civpers-overview-2025}.
  \item[\textbf{HRO focus.}] \ac{hro} primary functions are classification, staffing, employee relations,
  and labor relations; \ac{eeo} is not an \ac{hro} function~\autocite{edo-6-1-1-civpers-overview-2025}.
  \item[\textbf{EDQP link.}] \ac{edqp} requires completion of a local ``Supervision of Civilian Personnel''
  course, reinforcing the need to understand \ac{civpers} management basics~\autocite{edo-6-1-1-civpers-overview-2025}.
\end{description}

\subsection{Why CIVPERS matters}
\begin{description}
  \item[\textbf{Workforce scale.}] Civilian personnel dominate the manning at key technical commands and
  shipyards, making civilian leadership a core readiness driver for \ac{edo}s~\autocite{edo-6-1-1-civpers-overview-2025}.
\end{description}

\begin{longtblr}
  [
  caption={Illustrative civilian manning levels at selected activities},
  label={tab:civpers_manning_levels}
  ]{
  colspec={X[1.3]X[0.6]X[0.6]},
  rowhead=1
  }
  \toprule
  Activity & Military & Civilian \\
  \midrule
  NIWC LANT                     & 120 & 4{,}970 \\
  NIWC PAC                      & 200 & 4{,}800 \\
  Norfolk NSY                   & 810 & 10{,}600 \\
  NSWC Carderock                & 10  & 3{,}600 \\
  NSWC Port Hueneme             & 70  & 2{,}790 \\
  NSWC Panama City              & 45  & 1{,}410 \\
  Pearl Harbor NSY and IMF      & 490 & 5{,}270 \\
  Puget Sound NSY               & 35  & 12{,}000 \\
  Southwest Regional Maintenance Center & 750 & 1{,}150 \\
  \bottomrule
\end{longtblr}

\subsection{Myths and guidance sources}
\begin{description}
  \item[\textbf{Common myths.}] Misconceptions include believing civil servants are unaccountable,
  unions block mission success, and private industry rules always apply; leaders should verify facts
  and engage partners early~\autocite{edo-6-1-1-civpers-overview-2025}.
  \item[\textbf{Guidance stack.}] Key sources include Congress, the President, \ac{opm}, \ac{osd}, the
  Department of the Navy, \ac{navsea}, Human Resources Operation Center and Service Center channels,
  local \ac{hro} guidance, and negotiated labor agreements~\autocite{edo-6-1-1-civpers-overview-2025}.
\end{description}

\subsection{HRO primary functions}
\begin{description}
  \item[\textbf{Classification.}] Determine position titles and pay bands/grades~\autocite{edo-6-1-1-civpers-overview-2025}.
  \item[\textbf{Staffing.}] Fill vacancies and manage hiring actions~\autocite{edo-6-1-1-civpers-overview-2025}.
  \item[\textbf{Employee relations.}] Guide supervisors and employees on performance, discipline, and
  workplace issues~\autocite{edo-6-1-1-civpers-overview-2025}.
  \item[\textbf{Labor relations.}] Manage relationships and agreements with labor organizations~\autocite{edo-6-1-1-civpers-overview-2025}.
\end{description}

\subsection{Workforce shaping and classification}
\begin{description}
  \item[\textbf{Classification basics.}] Classification assigns a pay plan, occupational series, grade/band,
  and title; supervisors define duties while \ac{hro} applies \ac{opm} standards to determine the final
  classification~\autocite{edo-6-1-2-workforce-shaping-2025}.
  \item[\textbf{Pay systems.}] The most common systems are \ac{wg}, \ac{gs}, demonstration pay plans (e.g.,
  \ac{demo}/\ac{acqdemo}), \ac{sstm}, and \ac{ses}; understand grade/band structures before recruiting~\autocite{edo-6-1-2-workforce-shaping-2025}.
  \item[\textbf{Other pay and FLSA.}] Locality pay, premium pay, and \ac{flsa} exempt/non-exempt status
  depend on position duties; ensure the \ac{pd} matches actual responsibilities~\autocite{edo-6-1-2-workforce-shaping-2025}.
  \item[\textbf{Filling vacancies.}] Options include management reassignment, merit competition,
  competitive hiring, outsourcing, and mandatory placement programs; use \ac{dha} when authorized and
  plan timelines early~\autocite{edo-6-1-2-workforce-shaping-2025}.
  \item[\textbf{Downsizing levers.}] Workforce reductions rely on \ac{rif}, \ac{sip}, \ac{vera}, and furloughs;
  sequence actions carefully to preserve critical skills~\autocite{edo-6-1-2-workforce-shaping-2025}.
\end{description}

\subsection{Equal Employment Opportunity anchors}
\begin{description}
  \item[\textbf{Bases.}] The ten protected bases include race, color, religion, retaliation, sex (gender
  identity, sexual orientation, pregnancy), national origin, equal pay, age (40+), disability, and genetic
  information~\autocite{edo-6-1-3-eeo-2025}.
  \item[\textbf{Theories.}] Primary theories are disparate treatment, adverse impact, and denial of
  reasonable accommodation (disability or religion) unless undue hardship applies~\autocite{edo-6-1-3-eeo-2025}.
  \item[\textbf{Sexual harassment.}] Quid pro quo and hostile work environment are the two primary types;
  hostile environments require severe or pervasive conduct that interferes with work~\autocite{edo-6-1-3-eeo-2025}.
  \item[\textbf{Complaint process.}] The federal-sector process includes informal counseling, a formal
  complaint, and appeal; \ac{adr} may resolve issues earlier through mediation~\autocite{edo-6-1-3-eeo-2025}.
  \item[\textbf{Management responsibilities.}] Leaders must contact the \ac{eeo} office, respond quickly
  (including a 14-day inquiry window for sexual harassment), notify findings, and take appropriate
  action without reprisal~\autocite{edo-6-1-3-eeo-2025}.
\end{description}

\subsection{Employee relations essentials}
\begin{description}
  \item[\textbf{Supervisor and employee roles.}] Supervisors set expectations, reward good performance,
  and address poor conduct; employees are responsible for acceptable performance, attendance, and
  workplace behavior~\autocite{edo-6-1-4-employee-relations-2025}.
  \item[\textbf{Work schedules.}] Common schedules include basic, \ac{cws} (e.g., 5-4/9 or 4/10), and
  flexible schedules; overtime and telework are governed by local policy and \ac{flsa} status~\autocite{edo-6-1-4-employee-relations-2025}.
  \item[\textbf{Leave categories.}] Annual, sick, \ac{lwop}, \ac{awol}, \ac{fmla}, paid parental leave,
  military and court leave, the Voluntary Leave Transfer Program, and administrative excusal each
  carry specific approval and documentation rules~\autocite{edo-6-1-4-employee-relations-2025}.
  \item[\textbf{Performance management.}] \ac{dpmap} or \ac{acqdemo} processes require clear performance
  plans within 30 days, mid-term discussions, and \acp{pip} for unacceptable performance; probationary
  periods typically last one year~\autocite{edo-6-1-4-employee-relations-2025}.
  \item[\textbf{Awards and discipline.}] Incentive awards include monetary, time-off, and honorary
  recognition; discipline is progressive and must be timely, fair, and focused on correcting behavior~\autocite{edo-6-1-4-employee-relations-2025}.
\end{description}

\subsection{Within-grade increases and GS step timing}
\begin{description}
  \item[\textbf{Acceptable performance gates.}] \acp{wgi} require an \enquote{acceptable level of competence} determination (typically a \enquote{Fully Successful} rating); maintain clear standards and timely appraisals so eligibility decisions are defensible~\autocite{opm-wgi-factsheet-2024}.
\end{description}

\begin{longtblr}
  [
  caption={GS within-grade increase waiting periods},
  label={tab:wgi_waiting_periods}
  ]{
  colspec={X[1]X[1]},
  rowhead=1
  }
  \toprule
  Advancement to steps & Required waiting period at previous step \\
  \midrule
  Steps 2, 3, and 4    & 52 weeks (1 year)                        \\
  Steps 5, 6, and 7    & 104 weeks (2 years)                      \\
  Steps 8, 9, and 10   & 156 weeks (3 years)                      \\
  \bottomrule
\end{longtblr}

Assuming sustained acceptable performance, advancing from Step 1 to Step 10 takes 18 years (three years for Steps 1--4, six for Steps 4--7, and nine for Steps 7--10)~\autocite{opm-wgi-factsheet-2024,national-guard-technician-handbook-2023}. Exceptional performers may receive a \ac{qsi}, which moves them one step ahead of the normal schedule~\autocite{opm-wgi-factsheet-2024}.

\subsection{Performance vs.\ conduct pathways}
\begin{description}
  \item[\textbf{Supervisor baseline.}] Set and communicate standards, evaluate against them, and partner with \ac{hro} early; use coaching to reinforce desired behaviors~\autocite{edo-6-1-4-employee-relations-2025}.
\end{description}

\begin{longtblr}
  [
  caption={Differentiating performance and conduct issues},
  label={tab:performance_conduct_comparison}
  ]{
  colspec={X[1]X[1]X[1]},
  rowhead=1
  }
  \toprule
  Aspect             & Unacceptable performance                                       & Misconduct                                                                                                                          \\
  \midrule
  Definition         & Failure to perform assigned duties at an acceptable level      & Violation of a rule, law, or regulation (e.g., policy violations, absence without leave (AWOL), insubordination, security breaches) \\
  Focus of action    & Improve performance; not intended to punish                    & Correct behavior; penalties may be more severe depending on the offense                                                             \\
  Standard of proof  & Substantial evidence                                           & Preponderance of the evidence (higher standard than performance)                                                                    \\
  Guiding principles & Often involves a structured performance improvement plan (PIP) & Penalties must be just and reasonable; apply the Douglas Factors when selecting a penalty                                           \\
  \bottomrule
\end{longtblr}

\paragraph{Handling unacceptable performance.} Start with counseling and training, then use denial of a scheduled \ac{wgi} if performance remains below standard. A formal \ac{pip} (often about 90 days) documents required improvements; persistent failure can lead to change to lower grade or removal, with \ac{hro} coordination at each step~\autocite{edo-6-1-4-employee-relations-2025,national-guard-technician-handbook-2023}.

\paragraph{Addressing misconduct.} Apply progressive discipline calibrated to the offense: verbal or written counseling, reprimand, suspension, change to lower grade, or removal. The chosen penalty must be reasonable in light of the Douglas Factors (offense severity, job level, record, consistency, rehabilitation potential, and mission impact)~\autocite{edo-6-1-4-employee-relations-2025,national-guard-technician-handbook-2023}.

\paragraph{Employee rights and documentation.} Honor the Weingarten right to union representation during investigatory interviews that could lead to discipline; if representation is requested, either provide it, postpone, or end the interview~\autocite{edo-6-1-5-labor-relations-2025}. Keep actions non-discriminatory, coordinate on any \ac{eeo} complaints, and document counseling, training, and decisions to support subsequent actions~\autocite{edo-6-1-3-eeo-2025,edo-6-1-4-employee-relations-2025}.

\subsection{Labor relations touchpoints}
\begin{description}
  \item[\textbf{Definitions and coverage.}] Collective bargaining units are represented by an exclusive
  representative under a negotiated \ac{cba}; management officials, supervisors, confidential employees,
  certain HR and security roles, contractors, and uniformed members are excluded~\autocite{edo-6-1-5-labor-relations-2025}.
  \item[\textbf{Rights.}] Employees may join or refrain from unions, unions have a duty of fair
  representation, and management retains rights to mission, budget, organization, and assignments~\autocite{edo-6-1-5-labor-relations-2025}.
  \item[\textbf{Bargaining obligations.}] Agencies must bargain procedures and appropriate arrangements
  for management actions; impact and implementation bargaining requires notice, and unions must
  request bargaining or waive the right~\autocite{edo-6-1-5-labor-relations-2025}.
  \item[\textbf{Past practice.}] Established practices that are clear, consistent, long-standing, and mutually
  accepted can become enforceable; changes require notice-and-bargain steps~\autocite{edo-6-1-5-labor-relations-2025,FLRA-IRS-3-FLRA-513}.
  \item[\textbf{Formal discussions and Weingarten.}] Unions may attend formal discussions on conditions of
  employment and investigatory interviews when an employee requests representation~\autocite{edo-6-1-5-labor-relations-2025}.
  \item[\textbf{Official time and ULPs.}] Official time covers representational activities but not internal
  union business; \acp{ulp} are filed with \ac{flra} and must be reported within six months~\autocite{edo-6-1-5-labor-relations-2025}.
\end{description}

\subsection{CIVPERS wrap-up}
\begin{description}
  \item[\textbf{Core takeaways.}] Civilians are teammates with different experiences and expectations; do
  not assume familiarity with uniformed service norms, but do assume professionalism and commitment
  to the mission~\autocite{edo-6-1-6-civpers-wrap-up-2025}.
  \item[\textbf{Leader engagement.}] Invest in development, involve civilians in awards and performance
  processes, and communicate expectations consistently and transparently~\autocite{edo-6-1-6-civpers-wrap-up-2025}.
  \item[\textbf{Use your partners.}] Leverage \ac{hro} and senior civilians for guidance, and document
  decisions related to hiring, performance, and conduct~\autocite{edo-6-1-6-civpers-wrap-up-2025}.
  \item[\textbf{Practical application.}] Case studies reinforce how to apply rules on overtime, past
  practice, and workplace conduct while preserving fairness and mission focus~\autocite{edo-6-1-6-civpers-wrap-up-2025}.
\end{description}

%====================
% End of file
%====================
\IfSubfilesClassLoaded{
  \RenewDocumentCommand{\entryname}{}{\textbf{\color{Modern} Acronym}}
  \RenewDocumentCommand{\descriptionname}{}{\textbf{\color{Modern} Definition}}
  \printnoidxglossary[
    type=\acronymtype,
    title=Acronyms,
    style=long-booktabs
  ]}{}
\subsection{Coursebook cue: CIVPERS applied rules}

Coursebooks 6.1.2 through 6.1.6 test whether students can apply civilian
personnel rules to realistic supervisor scenarios, not just define the terms
\autocite{edo-6-1-2-workforce-shaping-2025,edo-6-1-3-eeo-2025,edo-6-1-4-employee-relations-2025,edo-6-1-5-labor-relations-2025,edo-6-1-6-civpers-wrap-up-2025}.

\begin{longtblr}[
  caption = {Vacancy-fill method selection cues.},
  label = {tab:vacancy-fill-cues},
]{
  width = \linewidth,
  colspec = {X[0.7,l] X[1.75,l]},
  rowhead = 1,
  hlines,
  vlines,
}
\textbf{Method} & \textbf{When it fits} \\
Competitive hiring & Best when the command needs a broad applicant pool or permanent placement. \\
Merit promotion & Best when filling from qualified federal employees under merit-system rules. \\
Reassignment or detail & Best for near-term mission need when the employee already has the required grade and qualification fit. \\
Temporary promotion & Best when higher-graded duties are temporary and properly classified. \\
Direct-hire or special authority & Best only when a valid authority applies and documentation supports its use. \\
\end{longtblr}

For FLSA, do not assume every salaried federal employee is exempt.  The
practical supervisor question is whether the position is exempt or non-exempt
and whether overtime was required or permitted beyond the basic work week
\autocite{edo-6-1-2-workforce-shaping-2025}.

\begin{longtblr}[
  caption = {EEO participant rights board cue.},
  label = {tab:eeo-participant-rights},
]{
  width = \linewidth,
  colspec = {X[0.7,l] X[1.8,l]},
  rowhead = 1,
  hlines,
  vlines,
}
\textbf{Participant} & \textbf{Supervisor cue} \\
Complainant & Has protected access to the EEO process and must not face reprisal for using it. \\
Responsible Management Official & Must cooperate, preserve records, avoid retaliation, and keep the response factual and professional. \\
Witness or advisor & Must be protected from interference or reprisal tied to participation. \\
\end{longtblr}

\begin{longtblr}[
  caption = {Leave-use scenario cues.},
  label = {tab:leave-use-cues},
]{
  width = \linewidth,
  colspec = {X[0.8,l] X[1.7,l]},
  rowhead = 1,
  hlines,
  vlines,
}
\textbf{Leave type} & \textbf{Supervisor cue} \\
Annual leave & Schedule subject to mission needs and applicable policy; avoid arbitrary denial. \\
Sick leave & Use for qualifying medical, family-care, bereavement, adoption, or exposure situations under governing rules. \\
Leave without pay & Discretionary unless an entitlement applies; document the basis and duration. \\
Family and medical leave & Check eligibility, covered reason, notice, certification, and job-protection rules. \\
Administrative leave & Use only when authorized; do not treat it as a general substitute for discipline or convenience. \\
\end{longtblr}

\begin{longtblr}[
  caption = {Labor-relations scenario cues.},
  label = {tab:labor-relations-scenario-cues},
]{
  width = \linewidth,
  colspec = {X[0.85,l] X[1.65,l]},
  rowhead = 1,
  hlines,
  vlines,
}
\textbf{Issue} & \textbf{Supervisor cue} \\
Official time & Some representational activity is statutory; other use depends on contract, bargaining, or negotiated procedure. \\
Formal discussion & Union notice and opportunity to attend may be required when conditions meet the statutory test. \\
Weingarten right & A bargaining-unit employee may request union representation during an investigatory examination that could lead to discipline. \\
ULP risk & Bypassing the union, refusing to bargain, retaliation, coercive statements, or unilateral changes can create unfair-labor-practice exposure. \\
\end{longtblr}

\paragraph{Board cue.}
For supervisor scenarios, identify the employee category, authority, required
process, documentation, employee or union right, and the behavior that would
create reprisal, EEO, FLSA, or ULP risk.

\end{document}
